% 外测度与 Carathéodory 扩张
\section{外测度}

\begin{wenxintishi}[阅读顺序]
外测度是「允许 ∞ 的长度」，只在可数并上满足次可加. 它比测度弱，因而容易构造；
后面再用 Carathéodory 条件把它裁成一个真正的测度.
\end{wenxintishi}

\begin{definition}[外测度]
设 $\mu^{*}:\mathcal{P}(X)\to[0,\infty]$ 满足 $\mu^{*}(\varnothing)=0$ 且对任意
$A\subseteq\bigcup_{n}E_n$ 有 $\mu^{*}(A)\le\sum_n\mu^{*}(E_n)$，则称 $\mu^{*}$
为 $X$ 上的外测度.
\end{definition}

\begin{lemma}[单调性]
外测度单调：若 $A\subseteq B$，则 $\mu^{*}(A)\le\mu^{*}(B)$.
\end{lemma}
\begin{proof}
$B\subseteq B\cup\varnothing\cup\varnothing\cup\cdots$，取 $E_1=B$、其余为空集即可.
\end{proof}

\section{Carathéodory 条件}

\begin{theorem}[Carathéodory 判别]
设 $\mu^{*}$ 为 $X$ 上的外测度，令
\[
\mathcal{M}=\{E\subseteq X:\forall\,A\subseteq X,\ \mu^{*}(A)=\mu^{*}(A\cap E)+\mu^{*}(A\setminus E)\}.
\]
则 $\mathcal{M}$ 是 $\sigma$-代数，且 $\mu^{*}\!|_{\mathcal{M}}$ 是测度.
\end{theorem}

\begin{proposition}[平移不变的外测度]
在 $\mathbb{R}$ 上由区间长度生成的外测度满足 $\mu^{*}(E+t)=\mu^{*}(E)$.
\end{proposition}

\begin{corollary}
Lebesgue 测度完备，且所有 Borel 集都可测.
\end{corollary}

\section{例题与练习}

\begin{example}[构造外测度]
设 $X=\mathbb{N}$，$\mu^{*}(A)=\sum_{n\in A}2^{-n}$. 验证 $\mu^{*}$ 是外测度，
并求出 $\mu^{*}(\mathbb{N})$.
\end{example}
\begin{solution}
空集对应零和；次可加性由非负项重排不变得到. 而
$\mu^{*}(\mathbb{N})=\sum_{n=1}^{\infty}2^{-n}=1$.
\end{solution}

\begin{Exercises}
  \item 证明有限集的外测度为零.
  \item 举出 $\mathbb{R}$ 的一个不可测子集存在的论证思路.
  \item 设 $\mu^{*}(E)<\infty$ 且 $E\in\mathcal{M}$，求 $\mu^{*}(E^{c})$.
\end{Exercises}
